\section{Experimental protocol}
\label{sec:protocol}

\subsection{Paired design}
\label{sec:paired}

A seed is not a random-number seed but \emph{the name of an experimental
condition}. For $\text{seed} = s$, every controller sees the same instance, the
same initial point, and the same minibatch order. The random stream used to
generate tasks is fully separated from the stream used to execute the optimizer,
so the instance is identical regardless of how much randomness a controller
consumes. This is a property of our runner, not a claim taken from the
literature; we verify it with the bitwise reproduction check of
Section~\ref{sec:identity}.

Statistics are computed on paired differences using the Wilcoxon signed-rank test
\citep{wilcoxon1945} and bootstrap confidence intervals for the median
\citep{efron1979bootstrap}. We write $p < 0.0001$ rather than $p = 0.0000$.

\subsection{Separation of seed roles}
\label{sec:seeds}

\begin{center}
\begin{tabular}{lll}
\toprule
Role & Seeds & Used for \\
\midrule
\texttt{CALIBRATION\_SEEDS} & 0, 1        & benchmark spec selection only \\
\texttt{SELECTION\_SEEDS}   & 2, 3, 4     & configuration selection only \\
\texttt{HELD\_OUT\_SEEDS}   & 100--109    & final effect estimation only \\
\bottomrule
\end{tabular}
\end{center}

\emph{The configuration was not chosen using held-out results.} It was fixed once
on development seeds in Section~\ref{sec:selection} and not reopened in
Section~\ref{sec:heldout}.

We disclose one remaining overlap. \texttt{quad\_d100\_k1e5} appears both in the
initial dev subset and in the final spec set, so
$(\texttt{quad\_d100\_k1e5}, \text{seed } 2)$ is the same instance in two phases.
This is one of the twelve instances used for configuration selection.

\subsection{Result identity}
\label{sec:identity}

Editing aggregation code or documentation must not cause the optimizer to rerun.
We separate identity into three layers:

\begin{center}
\begin{tabular}{ll}
\toprule
Identifier & Content \\
\midrule
\texttt{run\_semantics\_id} & only the optimizer settings this controller actually uses \\
\texttt{sweep\_id}          & the set of runs this execution requested \\
\texttt{aggregation\_id}    & the aggregation policy \\
\bottomrule
\end{tabular}
\end{center}

\texttt{git\_commit} and \texttt{code\_dirty} enter none of these identifiers and
are recorded as separate provenance. If a hash changed when only documentation
was edited, the meaning of ``which set was requested'' would be broken.

To check that this separation actually holds, we compared planner trajectories
before and after changes to identity, logging, and the baseline clock. On the same
CPU, single-threaded, with the same seed, the expectation is bitwise equality. All
$72$ predeclared pairs matched bitwise.

\subsection{Predeclared gates}
\label{sec:gates}

\begin{center}
\begin{tabular}{llr}
\toprule
Gate & Question & GO threshold \\
\midrule
\gate{A1} & Instantaneous headroom of the one-step absolute action space & $\geq 1.0$~nat \\
\gate{A2} & Is that headroom reachable with realistic multiplier actions? & $\geq 0.7$~nat \\
\gate{B}  & Is the action space the bottleneck? & $\geq 0.5$~nat \\
\gate{C1} & Does resetting the quota each step hurt? (diagnostic only) & $\geq 0.3$~nat \\
\gate{C2} & Is multi-step replanning better than one-step? & $\geq 0.3$~nat \\
\gate{C3} & Is state-observing replanning better than fixed execution? & $\geq 0.3$~nat \\
\gate{D}  & Cost-to-target headroom & $\geq 1.2\times$ \\
\bottomrule
\end{tabular}
\end{center}

A GO verdict on \gate{C2} requires \emph{both} conditions: an improvement must be
present, and plans of depth greater than one must actually be adopted. If only the
improvement is present while the depth stays at one, the contribution came from
search volume rather than from planning.

\gate{C1} uses a single-seed diagnostic baseline and is therefore not used for any
verdict.
